
\section{Preliminaries: scanning protocols, readout, and period data}
\label{sec:preliminaries}

\subsection{A controlled protocol class}
\label{sec:prelim_scanalg}

We work in a controlled class of auditable scan--readout protocols.
For concreteness, we focus on Kronecker-type scans on tori with algebraic (rational) kernels and explicit regularization rules.

\begin{definition}[Kronecker scan protocol]
\label{def:protocol}
A (Kronecker) scan protocol is a tuple
\[
\mathcal{P}=(\TT^d,\alpha,x_0,f,R),
\]
where $\TT^d=(\RR/\ZZ)^d$ is the $d$-torus, $\alpha\in\RR^d$ is a scan slope, $x_0\in\TT^d$ is an initial condition, $f$ is a readout kernel, and $R$ is a regularization rule producing a regulated kernel $f_R$.
The scan orbit is
\[
x_t = x_0+t\alpha \pmod 1,\qquad t\in\ZZ_{\ge 0}.
\]
\end{definition}

In this paper, the \emph{finite horizon} $N$ is treated as the primary resource variable.
We denote the orbit prefix by
\[
P_N=\{x_0,x_1,\dots,x_{N-1}\}\subset [0,1)^d.
\]
Additional resource parameters (e.g.\ truncation depth) are recorded as part of $R$ when relevant.

\subsection{Readout: finite-time averages}
\label{sec:prelim_readout}

\begin{definition}[Finite-horizon readout]
\label{def:readout}
Given a protocol $\mathcal{P}$ and horizon $N\ge 1$, the protocol readout is the empirical average
\[
\avgN{\mathcal{P}} := \frac{1}{N}\sum_{t=0}^{N-1} f_R(x_t).
\]
\end{definition}

The readout is the \emph{only} protocol output used in the closed layer: it is finite, computable, and reproducible given the protocol data.
Ensemble versions (averaging over multiple initial conditions) will be used in the numerical section.

\subsection{Period data and ideal values}
\label{sec:prelim_period_data}

To connect protocol outputs to classical constants, we treat the ideal target as a (regularized) period integral.

\begin{definition}[Period datum]
\label{def:period_datum}
A \emph{period datum} is a triple $W=(D,f,R)$ where $D\subset\RR^d$ is a measurable domain (in practice $D=[0,1)^d$), $f$ is a kernel, and $R$ is a regularization rule producing $f_R$.
Its (regularized) period value is
\[
\per(W) := \int_D f_R\,\dd x.
\]
\end{definition}

In the Kontsevich--Zagier sense, periods are values of integrals of algebraic differential forms over domains defined by algebraic inequalities \cite{KontsevichZagier2001}.
Our controlled protocol class is designed so that $\per(W)$ lies in, or is canonically related to, such a period class.

\subsection{Examples: basic period data used in this paper}
\label{sec:prelim_examples}

\begin{example}[$\log 2$]
On $D=[0,1]$, let $f(x)=\frac{1}{1+x}$ and $R=\mathrm{Id}$.
Then $W=(D,f,R)$ satisfies $\per(W)=\int_0^1 \frac{1}{1+x}\,\dd x=\log 2$.
\end{example}

\begin{example}[$\pi$]
On $D=[0,1]$, let $f(x)=\frac{4}{1+x^2}$ and $R=\mathrm{Id}$.
Then $W=(D,f,R)$ satisfies $\per(W)=\int_0^1 \frac{4}{1+x^2}\,\dd x=\pi$.
\end{example}

\begin{example}[$\pi^2$ and $\pi^3$ as cubical periods]
\label{ex:pi_powers}
Let $g(x)=\frac{4}{1+x^2}$ on $[0,1]$.
Define $f_2(x,y)=g(x)g(y)$ on $[0,1]^2$ and $f_3(x,y,z)=g(x)g(y)g(z)$ on $[0,1]^3$ (with $R=\mathrm{Id}$).
By Fubini's theorem,
\[
\int_{[0,1]^2} f_2\,\dd x\,\dd y=\left(\int_0^1 g(x)\,\dd x\right)^2=\pi^2,
\qquad
\int_{[0,1]^3} f_3\,\dd x\,\dd y\,\dd z=\left(\int_0^1 g(x)\,\dd x\right)^3=\pi^3.
\]
These product forms provide an explicit period realization of $\pi^2$ and $\pi^3$ within the same rational-kernel protocol class used throughout the paper.
\end{example}

\subsection{The holographic scan-to-period interface}
\label{sec:prelim_hsp}

The essential bridge is that a protocol determines a canonical period datum by forgetting orbit-level details while keeping the domain, kernel, and regularization.

\begin{definition}[Scan-to-period interface map]
\label{def:hsp}
For a protocol $\mathcal{P}=(\TT^d,\alpha,x_0,f,R)$ we define its associated period datum by
\[
\HSP(\mathcal{P}) := ([0,1)^d,f,R).
\]
When we restrict to the controlled class $\ScanAlg$, this assignment can be organized as a functor $\HSP:\ScanAlg\to\PerDatum$; in this paper we only need the object-level map and its invariance properties.
\end{definition}

\begin{remark}[Protocol equivalence and invariance]
\label{rem:protocol_equivalence}
Many changes in protocol presentation (e.g.\ torus translations, measure-preserving automorphisms, or equivalent regularization conventions) do not change $\per(\HSP(\mathcal{P}))$.
This invariance is the basis for treating period values as \emph{protocol invariants} rather than properties of a specific implementation.
\end{remark}

\begin{definition}[Affine protocol equivalence]
\label{def:protocol_equivalence}
Let $\mathcal{P}=(\TT^d,\alpha,x_0,f,R)$.
For $A\in \mathrm{GL}(d,\ZZ)$ with $|\det A|=1$ and $b\in\TT^d$, define the affine torus automorphism $\Phi(x)=Ax+b$.
We say that $\mathcal{P}'=(\TT^d,\alpha',x_0',f',R)$ is \emph{affine-equivalent} to $\mathcal{P}$ if
\[
\alpha'=A\alpha,\qquad x_0'=\Phi(x_0),\qquad f'_R = f_R\circ \Phi^{-1}
\]
(equivalently, $f'$ is the pullback of $f$ under $\Phi$ with the same regularization rule).
\end{definition}

\begin{lemma}[Period invariance under affine protocol equivalence]
\label{lem:period_invariance_affine}
If $\mathcal{P}'$ is affine-equivalent to $\mathcal{P}$ in the sense of Definition~\ref{def:protocol_equivalence}, then
\[
\per(\HSP(\mathcal{P}'))=\per(\HSP(\mathcal{P})).
\]
\end{lemma}

\begin{proof}
By definition, $\HSP(\mathcal{P})=([0,1)^d,f,R)$ and $\HSP(\mathcal{P}')=([0,1)^d,f',R)$ with $f'_R=f_R\circ \Phi^{-1}$.
Since $\Phi$ is measure-preserving on $\TT^d$ (it is an affine automorphism with $|\det A|=1$),
the change-of-variables formula yields
\[
\int_{[0,1)^d} f'_R\,\dd x=\int_{[0,1)^d} f_R\circ \Phi^{-1}\,\dd x=\int_{[0,1)^d} f_R\,\dd x.
\]
\end{proof}


